% !TEX program = pdflatex
\documentclass[3p,11pt]{elsarticle}
\usepackage{amsmath,amssymb}
\usepackage{array}
\usepackage{booktabs}
\usepackage[tableposition=top]{caption}
\usepackage{enumitem}
\usepackage{tikz}
\usepackage{pgfplots}
\pgfplotsset{compat=1.18}
\usepackage{url}
\urlstyle{same}
\usepackage[hidelinks]{hyperref}
\emergencystretch=3em
% first-page footer: the page number alone
\makeatletter
\def\ps@pprintTitle{\let\@oddhead\@empty\let\@evenhead\@empty
  \def\@oddfoot{\hfil\thepage\hfil}\let\@evenfoot\@oddfoot}
\makeatother

% ---- notation ---------------------------------------------------------------
\newcommand{\um}{\,$\mu$m}
\newcommand{\kms}{km\,s$^{-1}$}
\newcommand{\flux}{erg\,cm$^{-2}$\,s$^{-1}$}
\newcommand{\mA}{m\AA}
\newcommand{\sci}[2]{#1\times10^{#2}}
\newcommand{\ion}[2]{#1\,\textsc{#2}}
\newcommand{\nuc}[2]{$^{#1}$#2}
\newcommand{\lxl}{$L_X/L_\mathrm{bol}$}

\begin{document}

\begin{frontmatter}

\title{Dimming, orbits, clocks and abundances: archival tests for technosignatures among the stars}

\author{Charles C. Norton}

\begin{abstract}
A technology at a star could change the star's light, the motions of the bodies around
it, the timing of its variations or the composition of its surface, in ways that its
astrophysics does not. We test each kind of change against the archive that measures
it. In light, no star dims over a century, between Hipparcos and Gaia or within Gaia's
epochs without a natural cause, no transit deepens from Kepler to TESS in the light curves
themselves, the Kepler giants lack flux as often as they gain it,
no variable changes colour while holding its brightness or holds its brightness at
discrete levels, and no hot body is smaller than a white dwarf. In motion, no orbit is
too light for the light it shows, no eclipsing binary too
light for its eclipses, no transit faster than escape velocity, no planet pair in exact
commensurability, no orbit synchronised with its star beyond the tidal range, no
planet aimed at the Earth, and no planetary system or co-moving group arranged beyond
chance. In timing, no two stars share a clock beyond the period distribution, and
no star in the Earth's transit zone or on the ellipsoid of a nova flashes at the moment
it would see us or the nova. In composition, no dwarf carries fission products
or processed carbon and nitrogen without a natural cause. The one result that remains
open, a $5.5$ per cent polarisation assigned to an unreddened star in 1977, is the
subject of its own note.
\end{abstract}

\begin{keyword}
Technosignatures \sep Stars: variables: general \sep Binaries: general \sep
Stars: abundances \sep Astrometry
\end{keyword}

\end{frontmatter}

% =============================================================================
\section{Introduction}\label{sec:intro}
% =============================================================================

Most technosignature searches look for a signal; a technology also leaves the star it
occupies changed. Material in orbit dims and reddens the starlight it intercepts and
re-radiates it in the infrared~\cite{Dyson1960,Wright2014}, a structure placed to be seen
transits in a way a planet does not~\cite{Arnold2005}, and material deposited in a star
changes its surface composition~\cite{WhitmireWright1980}. The large surveys of the last
three decades measure every star near the Sun in brightness, colour, position, motion
and composition, at many epochs, and each of these changes has a limit the star's own
physics sets. This paper puts the limits to the surveys, with a null drawn from the same
data for each, and states what remains.

\subsection{Contributions}
\begin{enumerate}[leftmargin=*,itemsep=1pt]
\item Tests of starlight: dimming over a century, between Hipparcos and Gaia and within
Gaia's epochs, transit depths from Kepler to TESS, missing flux in asteroseismic giants, colour without brightness and
brightness at discrete levels, hot bodies smaller than white dwarfs, and dust at white
dwarfs.
\item Tests of motion: orbits too light for their light, eclipsers too light for their
eclipses, transits faster than escape velocity, exact commensurabilities, synchronous
orbits, aimed transits, the orientation of orbits and the clustering of co-moving stars.
\item Tests of timing and geometry: shared clocks, the Earth's transit zone at
opposition, the ellipsoids of novae and the stars in the beams of planetary radar.
\item Tests of composition: fission products and processed carbon and nitrogen in dwarfs.
\end{enumerate}

% =============================================================================
\section{Light}\label{sec:light}
% =============================================================================

\emph{Dimming over a century.} A structure built over a century would dim its star by a
measurable amount. In the Harvard plates of DASCH~\cite{Tang2013}, the $300$ nearest
sun-like stars of $G=8$ to $11$ have differential trends against comparison stars on the
same plates with a spread of $0.12$\,mag per century, and the two above $3.5\sigma$ have
slopes in their plate series that disagree with one common slope; injected fadings are
recovered in $7$, $60$ and $83$ per cent of cases at $0.4$, $0.6$ and $1.0$\,mag per
century.

\emph{Dimming between Hipparcos and Gaia.} For $15{,}189$ Hipparcos stars within
$100$\,pc fainter than $G=6$, Gaia's $G$ against the Hipparcos magnitude~\cite{Perryman1997}
after a colour term spreads by $0.0034$\,mag for constant single stars; $218$ stars are
fainter in Gaia by five robust deviations, every one with a natural cause: a Hipparcos
variability flag, a neighbour Hipparcos summed, a close companion that Gaia's $G$ misses,
a spotted rotator or an astrometric binary. Of the $531$ brighter in Gaia, $19$ have
none.

\emph{Dips within Gaia and Hipparcos epochs.} Of $130$ variables within $100$\,pc whose
Gaia epoch photometry dips with a skewness above $1.5$ outside the classifier's natural
classes, every dipper is a single epoch, a single transit with its neighbours undimmed,
a close binary measured differently at each scan angle, a periodic eclipse, reddening or
a dip in $G$ alone. Of $17{,}285$ Hipparcos stars, $23$ keep a dip seen in both fields of
view, the rate beyond $100$\,pc, and the deepest is an eclipsing binary in TESS. The
$28$ Gaia alerts within $2''$ of a star within $100$\,pc are flares on red and brown
dwarfs, two dwarf-nova outbursts, a young star's flare, a flare star's slow rise and a
distant dwarf nova behind a spurious parallax. Of the stars with both a Hipparcos-to-Gaia
deviation and a NEOWISE W2 trend, $671$, the highest combined value of optical dimming
and mid-infrared brightening arises by chance at its rank $99$ per cent of the time
against the W2 trends re-paired among stars of like brightness.

\emph{Transit depths from Kepler to TESS.} A structure under construction, or a swarm
filling in, deepens its transit at the same period and duration. Of $102$ Kepler objects
of interest that are TESS objects of interest within $6''$ and $0.2$ per cent in period,
three are deeper in TESS at $5\sigma$ with the duration held, and two shallower, the
null. Measured the same way in both missions' light curves, each transit's centre against
a line fitted to its sides, two of the deeper transits are deeper only in the
crowding-corrected flux, the pipeline's estimate of the neighbours' light in TESS's
$21''$ pixels, and one is deeper only in the TOI table; of the shallower, one is diluted
by light in the aperture beyond the crowding model and one is shallower only in the
table.

\emph{Missing bolometric flux.} A star's seismic radius and temperature give its
luminosity; a structure intercepting its light leaves Gaia's luminosity short of it.
Among $14{,}959$ single Kepler giants with seismic radii~\cite{Yu2018} and Gaia
luminosities~\cite{Berger2020}, deficits are as common as excesses at every threshold,
$31$ against $35$ beyond $4\sigma$, and the extreme deficits have a dwarf's luminosity
and a giant's oscillations: background giants in foreground dwarfs' apertures.

\emph{Colour without brightness, and brightness in levels.} A filter that changes a
star's colour while holding its brightness would be no spot or companion. Of $1073$ Gaia
variables brighter than $G=16$ whose $G$ scatter is below $0.003$\,mag while BP and RP
scatter by $0.02$ or more, $899$ exceed the quiet sources' tail; the $60$ highest, read
epoch by epoch, move BP and RP together, a neighbour entering both windows, or scatter in
one band alone, and none moves BP against RP. Conversely, of $1845$ ZTF spotted
variables~\cite{Chen2020}, nine have an amplitude ratio that is grey at $3\sigma$: eight
RS~CVn and BY~Dra systems above the main sequence, where the companion dilutes the spots'
colour, and a pre-main-sequence star in NGC~2264 occulted by its disk. In the TESS $20$\,s light curves
of the stars within $30$\,pc, $2096$ steps at $7\sigma$ with transitions of two minutes or
less are, $149$ of them, common to three or more stars of a sector, and the rest split
$1484$ up and $460$ down, flare rises; no light curve holds three levels each revisited.

\emph{Bodies too small or too blue.} A body hot by colour and smaller than a white dwarf
would need a source other than a star. Among $5278$ sources within $100$\,pc with
black-body colour temperatures above $7000$\,K, $22$ with clean astrometry and photometry
lie below $0.003\,R_\odot$, all of them catalogued white dwarfs, four in the list of
infrared-faint white dwarfs of Bergeron et al.~\cite{Bergeron2022}, whose collision-induced
absorption makes a $3000$\,K photosphere read hot by colour, and none with an ultraviolet
detection. In the Euclid Q1 catalogue, none of the $14$ point sources bluer than the
Rayleigh-Jeans slope in $I_E-H_E$ is so at $3\sigma$ with clean flags, unblended and
steady.

\emph{White dwarfs.} Dust at a white dwarf pollutes its atmosphere with the metals it
accretes, and warm material without pollution would be no disc of rock. Of $1434$ DESI
DR1 white dwarfs with unWISE W1 and W2 and parallaxes at $5\sigma$, $13$ exceed the
photosphere in both bands, five of them unblended with excesses of $900$ to $1500$\,K
and $1.6$ to $3.3$\,mag in W1, the excess of an unresolved brown dwarf rather than the
$0.3$ to $1$\,mag of a dust disc. Of $16{,}676$ white dwarfs within $100$\,pc~\cite{GentileFusillo2021},
$25$ have W3 or W4 above the photosphere with W1 and W2 photospheric; the strongest are an
extended source in the Galactic plane, an optical ghost, single-exposure detections and
cirrus.

% =============================================================================
\section{Motion}\label{sec:motion}
% =============================================================================

\emph{Light without mass.} A luminous object that weighs nothing moves its star's
photocentre and not its spectrum. Of $42$ Gaia DR3 astrometric orbits within $100$\,pc
with HARPS or HIRES velocities, nine have photocentre orbits the velocities exclude:
close and long-period companions, a period beside Gaia's scanning law, and blended
double-lined stars, two of which Gaia itself resolves on part of its transits.

\emph{Eclipsers too light for their eclipses.} An opaque screen eclipses without moving
its star. Among $956$ Gaia eclipsing binaries deeper than $0.75$\,mag with ten or more RVS
transits, the observed velocity amplitude over the least a main-sequence eclipser of the
depth's size demands has a median of $0.73$, and the ten below $0.1$ are a saturated
giant, an unphysical model depth, two binaries whose blended lines read low, a Kepler
system whose true period and depth differ from Gaia's, and five stars flat at Gaia's
period in ZTF and TESS: Gaia's depths and periods, not the velocities, fail.

\emph{Transits faster than escape velocity.} A bound body crosses its star no faster than
the star's surface escape speed, which sets a floor on a full transit's duration. Of the
Kepler threshold-crossing events with characterised hosts~\cite{Berger2020}, $281$ fall
below $0.8$ of the floor at their fitted geometry: $253$ on giants, $26$ pulsators, an
eclipsing binary, a KOI with centroid flags and a single cadence flagged for cosmic rays.
Of the TESS objects of interest, $242$ fall below the floor at zero impact parameter:
giants, contact binaries and pulsators, published false positives, and $42$ candidates
without an impact parameter in the catalogue. Their TESS light curves, folded at the
catalogue ephemeris and fitted with a trapezoid whose flat-bottom fraction bounds the
impact parameter, give $20$ a duration that a bound orbit reaches at some impact
parameter, four a V shape that leaves a grazing chord, and three a flat bottom whose
largest allowed impact parameter keeps the crossing within the escape speed; five show
no transit at the catalogue ephemeris, and ten have no TESS light curve in the public
archives.

\emph{Commensurability and synchrony.} An engineered system could hold its planets in
exact resonance or its orbits in step with its star. Of $2393$ planet pairs, none lies
within $10^{-5}$ of a small-integer ratio, and five within $10^{-4}$ against $3.4$, in the
librating chains of TOI-1136, Kepler-223 and HD~110067. Of $541$ Kepler planets at
periods of $10$\,d or more around hosts with rotation periods~\cite{McQuillan2013,Mazeh2015},
$25$ orbit within $2$ per cent of the rotation period, its double or its half, against
$23.4\pm4.6$ with the periods shuffled, and the four $1{:}1$ matches drift in rotational
phase by $0.17$ to $0.25$ of a cycle over four years, as spots do.

\emph{Aimed transits.} A structure placed to be seen from the Earth transits at a small
impact parameter. The fitted impact parameters of $2294$ Kepler candidates and planets
pile up at zero where the fits of injected planets do too, through poorly constrained
fits at the parameter's lower bound, and the $153$ planets with impact parameters
measured to $0.1$ hold $7$ below $0.05$ against $11.3$ expected.

\emph{Orbits and groups.} Of $1361$ Gaia astrometric companions that may be
substellar~\cite{Gaia2023nss}, $6.2$ per cent have an orbital pole within $15^\circ$ of a
Galactic pole against $6.3$ per cent of the stellar companions at the same latitudes;
the $74$ whose upper mass bound lies below $0.2\,M_\odot$ hold nine against $4.7\pm2.1$
($p=0.04$), one reading of a small set among several. Of
$6044$ APOGEE dwarfs within $100$\,pc with velocities, the $31$ groups within $10$\,pc and
$1$\,\kms\ are the Hyades and co-natal pairs, with metallicity spreads only at the
pipeline's temperature edge; and the $2977$ wide binaries of the coverage store share no
anomaly between their members beyond the permutations.

\emph{Imaging candidates.} Of the candidate companions of the Gemini/NICI and NaCo direct
imaging surveys with two or more epochs, the few off both the co-moving and the
background track are known brown dwarfs in orbital motion, a background star on its
track in Gaia DR3, and a faint source moving $200$\,mas\,yr$^{-1}$ of its own, as a field
brown dwarf would.

% =============================================================================
\section{Timing and geometry}\label{sec:timing}
% =============================================================================

\emph{Shared clocks.} Beacons built by one civilization could share a period. Among
$770{,}043$ ZTF periodic variables~\cite{Chen2020}, pairs more than $10''$ apart that agree
in period scale linearly with the tolerance, as a smooth period distribution gives, and
those that agree in phase as well number $42{,}112$ against $41{,}384$ to $42{,}020$ with
the phases shuffled; the largest groups sit on periods the catalogue repeats exactly.

\emph{The Earth's transit zone at opposition.} A star in the Earth's transit zone sees the
Earth transit the Sun at our opposition, the moment a signal timed to that event would
leave. Of the stars within $100$\,pc and $0.27^\circ$ of the ecliptic with TESS two-minute
data at opposition, $27$ windows of sectors $42$ to $46$ hold no lone flash against $0.37$
expected and no brightening against $1.12$, and $12$ windows of sectors $70$ to $72$ hold
one lone flash against $0.53$.

\emph{Nova ellipsoids.} A civilization that saw a nova could signal when its light reached
them, and stars on the ellipsoid with the nova and the Earth at its foci would be seen to
do so now~\cite{Davenport2022,Cabrales2024}. Of the stars within $100$\,pc on the ellipsoids of T~CrB's
1866 and 1946 eruptions, GK~Per, V603~Aql, RR~Pic, DQ~Her, CP~Pup, V1500~Cyg and the
supernovae of 1572 and 1604, observed by TESS at their arrival, the flashes and light
curves with flashes match the same stars' other sectors, the one excess carried by
active binaries whose flares follow the sectors.

\emph{Radar beams and beams between stars.} A star in the beam of a planetary radar
would reply after twice its light-travel time. The $813$ stars within $100$\,pc in the
beams of every published radar measurement from 1968 to 2025 are no closer than
$48.9$ light years, so no reply is due before 2117. And of the $2705$ stars within
$20$\,pc, $37$ pairs lie on opposite sides of the Earth within a $0.25^\circ$ beam, as many
as random directions give.

% =============================================================================
\section{Composition}\label{sec:composition}
% =============================================================================

\emph{Fission products.} Fission waste deposited in a star raises zirconium, molybdenum,
ruthenium, barium, neodymium and europium together in a pattern no nucleosynthesis
makes~\cite{WhitmireWright1980}. Among $50{,}742$ GALAH DR4 stars~\cite{Buder2025} with clean
abundances of the six, fission-like patterns are rarer than the mirror pattern at every
threshold, $71$ against $230$ above $5\sigma$, and the $23$ FGK dwarfs with molybdenum and
ruthenium all lie within $2.5\sigma$ of solar.

\emph{Processed carbon and nitrogen.} A dwarf whose surface carries the products of the
CNO cycle has gained them from a companion. Of $31{,}351$ APOGEE dwarfs, $91$ are
nitrogen-rich at $5\sigma$ against the locus in temperature and metallicity and none
carbon-rich; $57$ of the $91$ are Gaia binaries against $28$ per cent of random dwarfs, the
mark of transfer from an evolved companion, $78$ carry binary, s-process, sodium,
rotation, evolution or temperature-edge flags, and the $13$ that carry none all carry
the pipeline's microturbulence grid-edge warning, which biases carbon and nitrogen in
dwarfs.

\begin{table}[t]
\centering\small
\caption{The tests, their samples and results.}
\label{tab:tests}
\begin{tabular}{>{\raggedright\arraybackslash}p{0.24\textwidth} >{\raggedright\arraybackslash}p{0.3\textwidth} >{\raggedright\arraybackslash}p{0.38\textwidth}}
\toprule
Test & Sample & Result \\
\midrule
Century dimming & $300$ DASCH stars & no consistent trend \\
Hipparcos to Gaia & $15{,}189$ stars & every dimming natural \\
Gaia and Hipparcos dips & $130$; $17{,}285$ stars & natural, at the background rate \\
Transit depths & $102$ KOI-TOI pairs & crowding corrections, catalogue depths \\
Missing flux & $14{,}959$ giants & deficits as common as excesses \\
Colour without brightness & $1073$ Gaia variables & neighbours in the windows \\
Discrete levels & $1806$ TESS sectors & flare rises and shared steps \\
Too small or too blue & $5278$; $14$ sources & infrared-faint white dwarfs; flags \\
White-dwarf dust & $1434$; $16{,}676$ white dwarfs & brown dwarfs, artefacts \\
Light without mass & $42$ orbits & companions and blends \\
Light eclipsers & $956$ binaries & Gaia's depths and periods \\
Escape velocity & Kepler, TESS events & giants, pulsators, false positives, chords \\
Commensurability & $2393$ pairs & librating chains \\
Synchrony & $541$ planets & at the chance rate \\
Aimed transits & $2294$ planets & the fits' pile-up \\
Orbit poles, groups & $1361$ orbits; $6044$ dwarfs & at chance; co-natal \\
Shared clocks & $770{,}043$ variables & at the period distribution \\
Transit zone at opposition & $39$ windows & at the background rate \\
Nova ellipsoids & ten eruptions & at the background rate \\
Fission products & $50{,}742$ stars & rarer than the mirror \\
Processed CNO & $31{,}351$ dwarfs & binaries and pipeline flags \\
\bottomrule
\end{tabular}
\end{table}

% =============================================================================
\section{Related work}\label{sec:related}
% =============================================================================

Dyson~\cite{Dyson1960} proposed the infrared signature of a star's surroundings
reprocessed, and Wright et al.~\cite{Wright2014} set out the search for it; Arnold~\cite{Arnold2005}
proposed artificial transits, Whitmire and Wright~\cite{WhitmireWright1980} fission waste in
a star, and Davenport et al.~\cite{Davenport2022,Cabrales2024} searched the ellipsoid of
SN~1987A for signals timed to it. The tests here rest on the surveys' own catalogues: Hipparcos and
Gaia~\cite{Perryman1997,Gaia2023nss}, DASCH~\cite{Tang2013}, the Kepler giants'
asteroseismology~\cite{Yu2018} and luminosities~\cite{Berger2020}, ZTF's periodic
variables~\cite{Chen2020}, the Kepler rotation periods~\cite{McQuillan2013,Mazeh2015}, the
Gaia white dwarfs~\cite{GentileFusillo2021,Bergeron2022} and GALAH~\cite{Buder2025}.

% =============================================================================
\section{Conclusion}\label{sec:conclusion}
% =============================================================================

No test finds a star changed in light, motion, timing or composition beyond what its
astrophysics and the surveys' own systematics give. Dimming is variability, blending or
an epoch's error; deepened transits are crowding corrections and catalogue depths; missing
flux is a background giant; colour without brightness is a
neighbour; the hot small bodies are infrared-faint white dwarfs; light orbits are blends
and companions; light eclipsers are Gaia's depths; fast transits are giants, pulsators
and chords away from the centre; commensurabilities are librating chains; aimed transits are the fits'
degeneracy; shared clocks are the period distribution; the transit zone and the nova
ellipsoids hold flashes at the background rate; and the abundance outliers are binaries
and pipeline flags. The reach of each test is its survey's: the century test reaches
fadings of about $0.6$\,mag per century, the Hipparcos-to-Gaia test $0.02$\,mag over
$25$ years, the missing-flux test deficits of $45$ per cent, and the motion tests the
precision of Gaia's astrometry and the archives' velocities.

\section*{Acknowledgements}

This work uses Gaia DR3 and Hipparcos (ESA; Gaia DPAC), DASCH (Harvard College
Observatory), Kepler and TESS through MAST, ZTF through IRSA, NEOWISE and AllWISE (IRSA),
APOGEE (SDSS-IV), GALAH DR4, DESI DR1 through the Astro Data Lab, Euclid Q1 (ESA), the
NASA Exoplanet Archive and ExoFOP, the Gemini/NICI and VLT/NaCo survey tables, JPL's radar
astrometry, and SIMBAD and VizieR at CDS, Strasbourg.

\begin{thebibliography}{99}\small\raggedright
\setlength{\itemsep}{1pt}

\bibitem{Dyson1960}
F.~J. Dyson, Search for artificial stellar sources of infrared radiation, \emph{Science}
\textbf{131}, 1667 (1960).

\bibitem{Wright2014}
J.~T. Wright, B.~Mullan, S.~Sigurdsson and M.~S. Povich, The $\hat{G}$ infrared search for
extraterrestrial civilizations with large energy supplies. I. Background and
justification, \emph{Astrophys. J.} \textbf{792}, 26 (2014).

\bibitem{Arnold2005}
L.~F.~A. Arnold, Transit light-curve signatures of artificial objects, \emph{Astrophys.
J.} \textbf{627}, 534 (2005).

\bibitem{WhitmireWright1980}
D.~P. Whitmire and D.~P. Wright, Nuclear waste spectrum as evidence of technological
extraterrestrial civilizations, \emph{Icarus} \textbf{42}, 149 (1980).

\bibitem{Tang2013}
S.~Tang, J.~Grindlay, E.~Los and M.~Servillat, Improved photometry for the DASCH pipeline,
\emph{Publ. Astron. Soc. Pac.} \textbf{125}, 857 (2013).

\bibitem{Perryman1997}
M.~A.~C. Perryman, L.~Lindegren, J.~Kovalevsky, et al., The HIPPARCOS catalogue,
\emph{Astron. Astrophys.} \textbf{323}, L49 (1997).

\bibitem{Yu2018}
J.~Yu, D.~Huber, T.~R. Bedding, et al., Asteroseismology of 16,000 Kepler red giants:
global oscillation parameters, masses, and radii, \emph{Astrophys. J. Suppl.}
\textbf{236}, 42 (2018).

\bibitem{Berger2020}
T.~A. Berger, D.~Huber, J.~L. van Saders, et al., The Gaia-Kepler stellar properties
catalog. I. Homogeneous fundamental properties for 186,301 Kepler stars, \emph{Astron.
J.} \textbf{159}, 280 (2020).

\bibitem{Chen2020}
X.~Chen, S.~Wang, L.~Deng, R.~de~Grijs, M.~Yang and H.~Tian, The Zwicky Transient Facility
catalog of periodic variable stars, \emph{Astrophys. J. Suppl.} \textbf{249}, 18 (2020).

\bibitem{Bergeron2022}
P.~Bergeron, P.~Dufour, G.~Fontaine, et al., On the nature of ultracool white dwarfs:
not so cool after all, \emph{Astrophys. J.} \textbf{934}, 36 (2022).

\bibitem{GentileFusillo2021}
N.~P. Gentile Fusillo, P.-E. Tremblay, E.~Cukanovaite, et al., A catalogue of white dwarfs
in Gaia EDR3, \emph{Mon. Not. R. Astron. Soc.} \textbf{508}, 3877 (2021).

\bibitem{McQuillan2013}
A.~McQuillan, T.~Mazeh and S.~Aigrain, Stellar rotation periods of the Kepler objects of
interest: a dearth of close-in planets around fast rotators, \emph{Astrophys. J. Lett.}
\textbf{775}, L11 (2013).

\bibitem{Mazeh2015}
T.~Mazeh, H.~B. Perets, A.~McQuillan and E.~S. Goldstein, Photometric amplitude
distribution of stellar rotation of KOIs: indication for spin-orbit alignment of cool
stars and high obliquity for hot stars, \emph{Astrophys. J.} \textbf{801}, 3 (2015).

\bibitem{Gaia2023nss}
Gaia Collaboration, F.~Arenou, C.~Babusiaux, et al., Gaia Data Release 3: stellar
multiplicity, a teaser for the hidden treasure, \emph{Astron. Astrophys.} \textbf{674},
A34 (2023).

\bibitem{Davenport2022}
J.~R.~A. Davenport, B.~Cabrales, S.~Sheikh, et al., Searching the SETI ellipsoid with
Gaia, \emph{Astron. J.} \textbf{164}, 117 (2022).

\bibitem{Cabrales2024}
B.~Cabrales, J.~R.~A. Davenport, S.~Z. Sheikh, et al., Searching the SN 1987A SETI
ellipsoid with TESS, arXiv:2402.11037 (2024).

\bibitem{Buder2025}
S.~Buder, J.~Kos, E.~X. Wang, et al., The GALAH survey: data release 4,
arXiv:2409.19858 (2024).

\end{thebibliography}

\end{document}
